%%%%%%%%%%%%%%%%%%%%%%%%%%%%%%%%%%%%%%%%%
% CSAS 2026 Poster — EPAA: A Two-Stage Approach to
% Quantifying Scoring Talent in the NBA
%
% Based on the Jacobs Landscape Poster LaTeX Template
% Version 1.0 (29/03/13)
% License: CC BY-NC-SA 3.0
%%%%%%%%%%%%%%%%%%%%%%%%%%%%%%%%%%%%%%%%%

%----------------------------------------------------------------------------------------
% PACKAGES AND OTHER DOCUMENT CONFIGURATIONS
%----------------------------------------------------------------------------------------

\documentclass[final]{beamer}

\usepackage[scale=1.24]{beamerposter}

\usetheme{confposter}

\setbeamercolor{block title}{fg=ngreen,bg=white}
\setbeamercolor{block body}{fg=black,bg=white}
\setbeamercolor{block alerted title}{fg=white,bg=dblue!70}
\setbeamercolor{block alerted body}{fg=black,bg=dblue!10}

%-----------------------------------------------------------
% Define the column widths and overall poster size
% 3-column layout: sepwid = 0.024, onecolwid = 0.3, twocolwid = 0.624
%-----------------------------------------------------------

\newlength{\sepwid}
\newlength{\onecolwid}
\newlength{\twocolwid}
\newlength{\threecolwid}

\setlength{\paperwidth}{48in}
\setlength{\paperheight}{50in}
%\setlength{\paperheight}{36in}

\setlength{\sepwid}{0.024\paperwidth}
\setlength{\onecolwid}{0.3\paperwidth}
\setlength{\twocolwid}{0.624\paperwidth}
\setlength{\threecolwid}{0.948\paperwidth}
\setlength{\topmargin}{-0.5in}

%-----------------------------------------------------------

\usepackage{graphicx}
\usepackage{booktabs}
\usepackage{amsmath}
\usepackage{amssymb}
\usepackage{hyperref}
\usepackage{qrcode}  % For QR code generation

%----------------------------------------------------------------------------------------
% TITLE SECTION
%----------------------------------------------------------------------------------------

% \title{A Two-Stage Approach to Quantifying Scoring Talent in the NBA}
% 
% \author{Alex Susi}
% 
% \institute{CSAS 2026 --- Poster \#20}

%----------------------------------------------------------------------------------------
\setbeamertemplate{headline}{}
\newsavebox{\titlebox}
\begin{document}

\addtobeamertemplate{block end}{}{\vspace*{2ex}}
\addtobeamertemplate{block alerted end}{}{\vspace*{2ex}}

\setlength{\belowcaptionskip}{2ex}
\setlength\belowdisplayshortskip{2ex}

% \begin{frame}[t]
% 
% \begin{columns}[T]
%   \begin{column}{0.15\paperwidth}
%     \centering
%     
\includegraphics[height=0.4\linewidth]{gaTech.png}
%   \end{column}
%   \begin{column}{0.7\paperwidth}
%     \centering
%     {\Huge\textbf{A Two-Stage Approach to Quantifying NBA Scoring Talent}}\\[2ex]
%     {\LARGE Alex Susi}\\[1ex]
%     {\LARGE CSAS 2026 ---  Georgia Institute of Technology}
%   \end{column}
%   \begin{column}{0.15\paperwidth}
%     \centering
%     \includegraphics[height=0.4\linewidth]{ball.png}
%   \end{column}
% \end{columns}

\begin{frame}[t]

\savebox{\titlebox}{%
  \begin{minipage}{0.7\paperwidth}
    \centering
    {\Huge\textbf{A Two-Stage Approach to Quantifying NBA Scoring Talent}}\\[2ex]
    {\LARGE Alex Susi}\\[1ex]
    {\LARGE CSAS 2026 --- Georgia Institute of Technology}
  \end{minipage}%
}

\begin{columns}[T]
  \begin{column}{0.15\paperwidth}
    \centering
    
\includegraphics[height=\ht\titlebox]{gaTech.png}
  \end{column}
  \begin{column}{0.7\paperwidth}
    \usebox{\titlebox}
  \end{column}
  \begin{column}{0.15\paperwidth}
    \centering
    \includegraphics[height=\ht\titlebox]{ball.png}
  \end{column}
\end{columns}
\vspace{0.5cm}
\begin{center}
\color{dblue}\rule{0.95\paperwidth}{0.15cm}
\end{center}
\vspace{0.5cm}

\begin{columns}[t]

\begin{column}{\sepwid}\end{column}


%========================================================================================
% COLUMN 1
%========================================================================================

\begin{column}{\onecolwid}

%----------------------------------------------------------------------------------------
% INTRODUCTION
%----------------------------------------------------------------------------------------

\begin{alertblock}{Introduction}

The ultimate objective of basektball is to score more points than the opponent. Consequentially, a fundamental challenge in basketball analytics is determining who the best scorers are. However, not all shot attempts are created equal. A wide-open dunk in transition, an open catch-and-shoot corner three-pointer, and a contested midrange stepback all have different expected conversion rates. Treating them identically conflates \textbf{shot quality} with \textbf{shooting skill}.

\vspace{2ex}

%\begin{center}
%\begin{tabular}{cc}
%\includegraphics[width=0.42\linewidth]{img_dunk.jpg} &
%\includegraphics[width=0.42\linewidth]{img_contested3.jpg} \\
%{\small Wide-open dunk} & {\small Contested deep three}
%\end{tabular}
%\end{center}

%\vspace{0.5ex}

Despite their importance for player evaluation, common public metrics such as FG\%, eFG\%, and even TS\% are heavily confounded by shot selection and defensive assignment. These are simple summaries of shooting outcomes that are much more context-dependent than implied by a single percentage.

\vspace{2ex}

We introduce \textbf{Expected Points Above Average (EPAA)}, a two-stage framework that:

\begin{enumerate}
\item Estimates \textit{shot quality} using contextual features.
\item Estimates \textit{shooting talent} via Bayesian hierarchical modeling.
\end{enumerate}

\vspace{2ex}

EPAA produces \textbf{league-adjusted metrics} for shooters, individual defenders, and team defense, such that 0 is league-average. This work uses public NBA tracking data from approximately 200,000 field goal attempts during the 2024–25 regular season, accessed via the \texttt{hoopR} R package.

\end{alertblock}

%----------------------------------------------------------------------------------------
% DATA
%----------------------------------------------------------------------------------------

% \begin{block}{Data}
% 
% \textbf{Source:} NBA Stats API accessed via the \texttt{hoopR} R package.
% 
% \vspace{0.5ex}
% 
% \textbf{Coverage:} ${\sim}200{,}000$ field goal attempts from the 2024--25 NBA regular season.
% 
% \vspace{0.5ex}
% 
% \textbf{Feature engineering pipeline} (\texttt{01\_data\_pipeline.R}):
% \begin{itemize}
% \item Play-by-play data with lineup tracking and shot clock reconstruction
% \item \textbf{Defender imputation:} Closest defender on each shot reconstructed from NBA matchup data, lineup tracking, and period-level co-occurrence weighting
% \item Fatigue features (minutes in game and stint)
% \item Shooting foul linkage for future True EPAA extension
% \end{itemize}
% 
% \vspace{0.5ex}
% 
% Shots are partitioned into three \textbf{shot families}: \textit{rim} (dunks, layups, tip-ins), \textit{mid-range} (j2), and \textit{three-point} (j3). Separate models are fit for each family to capture fundamentally different shooting mechanics and defensive dynamics.
% 
% \end{block}

%----------------------------------------------------------------------------------------
% METHODOLOGY
%----------------------------------------------------------------------------------------

% \begin{block}{Methodology}
% 
% The EPAA framework uses a \textbf{two-stage pipeline} that cleanly separates \textit{what happened} from \textit{who did it}:
% 
% \vspace{1ex}
% 
% \begin{center}
% \textbf{Stage 1: Shot Difficulty (XGBoost)} \\[0.5ex]
% Context-only features $\xrightarrow{\hspace{1cm}}$ Expected FG\% (xFG) \\[0.3ex]
% {\small\textit{``How hard was this shot for a league-average player?''}}
% \end{center}
% 
% \vspace{0.5ex}
% $\Downarrow$ \hfill {\small xFG passed as input to Stage 2}
% 
% \vspace{0.5ex}
% 
% \begin{center}
% \textbf{Stage 2: Talent Estimation (Bayesian Stan)} \\[0.5ex]
% xFG + make/miss $\xrightarrow{\hspace{1cm}}$ Player \& defender skill \\[0.3ex]
% {\small\textit{``Who is good, after controlling for shot difficulty?''}}
% \end{center}
% 
% \vspace{1ex}
% 
% \textbf{Key design principle:} Player and defender \textit{identity} is intentionally excluded from Stage 1. XGBoost sees only the shot situation (distance, angle, defender proximity, shot clock, etc.). This ensures Stage 2 random effects capture genuine skill, not confounded context.
% 
% \end{block}


%----------------------------------------------------------------------------------------



%----------------------------------------------------------------------------------------
% METHODOLOGY
%----------------------------------------------------------------------------------------

\begin{block}{Methodology}

\textbf{Stage 1: Shot Quality} \\[0.5ex]
\begin{itemize}
\item Expected Field Goal Percentage (xFG) is estimated using three separate XGBoost classifiers (one per shot type: rim, mid-range, three-point), trained on context features: shot distance, shot angle, defender proximity, dribble count, shot clock, player fatigue, and play type indicators.
\item Player and defender identities are intentionally excluded in these models, so they describe: \textit{"How hard was this shot for a league-average player guarded by a league-average defender?"}
\end{itemize}

\vspace{2ex}

\textbf{Stage 2: Shooting Talent} \\[0.5ex]
\begin{itemize}
\item Residual variation of shot difficulty is modeled using a Bayesian hierarchical model in Stan with random effects for shooters, individual defenders, and defensive teams.
\item Partial pooling shrinks low-volume players toward the league mean while preserving separation for high-volume stars, producing skill parameters with uncertainty quantification rather than simple point estimates.
\end{itemize}




\end{block}








%----------------------------------------------------------------------------------------

\end{column}

\begin{column}{\sepwid}\end{column}

%========================================================================================
% COLUMN 2
%========================================================================================

\begin{column}{\onecolwid}

%----------------------------------------------------------------------------------------
% SHOT QUALITY MODELS (STAGE 1)
%----------------------------------------------------------------------------------------

\begin{block}{Stage 1 Results: Shot Quality Models}

% Three separate XGBoost classifiers are trained on \textbf{context-only features} with an 85/15 temporal train/test split:
% 
% \vspace{0.5ex}
% 
% \begin{table}
% \begin{tabular}{l l}
% \toprule
% \textbf{Feature} & \textbf{Description} \\
% \midrule
% \texttt{shot\_distance} & Feet from basket \\
% \texttt{shot\_angle} & Degrees from baseline center \\
% \texttt{defender\_dist\_level} & Categorical closeness of defender \\
% \texttt{dribble\_level} & Number of dribbles before shot \\
% \texttt{shot\_clock} & Seconds remaining on shot clock \\
% \texttt{player\_min\_game/stint} & Fatigue proxies \\
% \texttt{is\_dunk} & Dunk indicator (rim only) \\
% \texttt{is\_fastbreak, is\_2ndchance} & Play-type indicators \\
% \bottomrule
% \end{tabular}
% \end{table}
% 
% \vspace{0.5ex}

Feature importance varies substantially across shot families, as revealed by SHAP analysis. This validated the decision to fit separate models per shot type rather than a single pooled classifier.

\vspace{0.5ex}

\begin{figure}
\includegraphics[width=\linewidth]{shap_rim.png}
\caption{Rim: Dunks and defender distance dominate}
\end{figure}

\begin{figure}
\includegraphics[width=\linewidth]{shap_j2.png}
\caption{Mid-range: Defender distance is the primary driver}
\end{figure}

\begin{figure}
\includegraphics[width=\linewidth]{shap_j3.png}
\caption{Three-point: Shot angle and distance most important}
\end{figure}

\end{block}

%----------------------------------------------------------------------------------------

\end{column}

\begin{column}{\sepwid}\end{column}

%========================================================================================
% COLUMN 3
%========================================================================================

\begin{column}{\onecolwid}

%----------------------------------------------------------------------------------------
% EPAA MODEL (STAGE 2)
%----------------------------------------------------------------------------------------

% \begin{block}{Stage 2: EPAA Bayesian Hierarchical Model}
% 
% {\small
% \textbf{Likelihood:}
% \begin{aligned}
%   y_n &\sim \mathrm{Bernoulli}\!\bigl(\mathrm{logit}^{-1}(\eta_n)\bigr),
%     \quad n = 1, \dots, N 
% \end{aligned}
% 
% \textbf{Linear predictor:}
% \begin{aligned}
%   \eta_n &= \underbrace{\alpha_{t_n}
%     + \beta_{t_n} \cdot \mathrm{xfg\_logit}_n
%     }_{\text{calibrated Stage 1}}
%     + \underbrace{a^{\mathrm{player}}_{j[n],\,t_n}
%     }_{\text{shooter skill}}
%     + \underbrace{a^{\mathrm{def}}_{d[n],\,t_n}
%     }_{\text{indiv.\ defender}}
%     + \underbrace{a^{\mathrm{team}}_{k[n],\,t_n}
%     }_{\text{def.\ scheme}} 
% \end{aligned}
% 
% \textbf{Random effects} (non-centered parameterization):
% \begin{aligned}
%   a^{\mathrm{player}}_{j,t} &= \mu^{\mathrm{vol}}_t
%     \cdot x^{\mathrm{vol}}_{j,t}
%     + \sigma^{\mathrm{player}}_t \cdot z^{\mathrm{player}}_{j,t},
%     &z^{\mathrm{player}}_{j,t}
%     &\sim \mathcal{N}(0, 1)  \\[4pt]
%   a^{\mathrm{def}}_{d,t}
%     &= \sigma^{\mathrm{def}}_t
%     \cdot z^{\mathrm{def}}_{d,t},
%     &z^{\mathrm{def}}_{d,t}
%     &\sim \mathcal{N}(0, 1)  \\[4pt]
%   a^{\mathrm{team}}_{k,t}
%     &= \sigma^{\mathrm{team}}_t
%     \cdot z^{\mathrm{team}}_{k,t},
%     &z^{\mathrm{team}}_{k,t}
%     &\sim \mathcal{N}(0, 1) 
% \end{aligned}
% 
% \textbf{Priors:}
% \begin{aligned}
%   \alpha_t &\sim \mathcal{N}(0,\; 0.25^2)  \\
%   \beta_t  &\sim \mathcal{N}^{+}(1,\; 0.25^2)  \\[4pt]
%   \mu^{\mathrm{vol}}_t &\sim \mathcal{N}(0,\; 0.3^2)  \\[4pt]
%   \sigma^{\mathrm{player}}_t
%     &\sim \mathcal{N}^{+}\!\bigl(0,\; s^{\mathrm{player}}_t\bigr),
%     \quad s^{\mathrm{player}}
%     = (0.30,\; 0.30,\; 0.25)  \\
%   \sigma^{\mathrm{def}}_t
%     &\sim \mathcal{N}^{+}(0,\; 0.20)
%     \quad \forall\; t  \\
%   \sigma^{\mathrm{team}}_t
%     &\sim \mathcal{N}^{+}\!\bigl(0,\; s^{\mathrm{team}}_t\bigr),
%     \quad s^{\mathrm{team}}
%     = (0.12,\; 0.10,\; 0.06) 
% \end{aligned}
% }
% 
% \vspace{0.3ex}
% 
% {\small Shot family index $t \in \{1, 2, 3\}$ for rim, mid-range, and three-point. Family-specific prior scales encode basketball structure: defensive team scheme effects are largest at the rim and smallest on threes. The model uses \texttt{reduce\_sum} parallelization for efficient MCMC sampling in Stan.}
% 
% \end{block}



%----------------------------------------------------------------------------------------
% EPAA MODEL (STAGE 2)
%----------------------------------------------------------------------------------------

\begin{block}{Stage 2: EPAA Bayesian Hierarchical Model}

For each shot attempt $n = 1, \dots, N$, let:
\[
\begin{aligned}
  y_n &\in \{0,1\} & &\text{shot outcome (make/miss)} \\
  t_n &\in \{1,2,3\} & &\text{shot type (rim, mid-range, three-point)}
\end{aligned}
\]

\vspace{1ex}

\textbf{Likelihood and Linear Predictor}
\[
\begin{aligned}
  y_n \sim \mathrm{Bernoulli}\!\bigl(\mathrm{logit}^{-1}(\eta_n)\bigr) \\
  \eta_n = \underbrace{\alpha_{t_n}
    + \beta_{t_n} \cdot \mathrm{xfg\_logit}_n
    }_{\text{Stage\,1 prediction}}
    + \underbrace{a^{\mathrm{p}}_{j[n],\,t_n}
    }_{\text{shooter skill}}
    + \underbrace{a^{\mathrm{f}}_{d[n],\,t_n}
    }_{\substack{\text{individual} \\ \text{defender}}}
    + \underbrace{a^{\mathrm{m}}_{k[n],\,t_n}
    }_{\substack{\text{team} \\ \text{defense}}}
\end{aligned}
\]

The intercept $\alpha_t$ and slope $\beta_t$ recalibrate the Stage~1 output per shot type. Three sets of \textbf{random effects} decompose residual skill.

\vspace{1ex}

\textbf{Random Effects}
\[
\begin{aligned}
  a^{\mathrm{p}}_{j,t} &= \mu^{\mathrm{vol}}_t
    \cdot x^{\mathrm{vol}}_{j,t}
    + \sigma^{\mathrm{p}}_t \cdot z^{\mathrm{p}}_{j,t},
    &z^{\mathrm{p}}_{j,t}
    &\sim \mathcal{N}(0,\, 1) \text{ for all players } j = 1,\dots,J \\[4pt]
  a^{\mathrm{f}}_{d,t}
    &= \sigma^{\mathrm{f}}_t
    \cdot z^{\mathrm{f}}_{d,t},
    &z^{\mathrm{f}}_{d,t}
    &\sim \mathcal{N}(0,\, 1) \text{ for all defenders } d = 1,\dots,D  \\[4pt]
  a^{\mathrm{m}}_{k,t}
    &= \sigma^{\mathrm{m}}_t
    \cdot z^{\mathrm{m}}_{k,t},
    &z^{\mathrm{m}}_{k,t}
    &\sim \mathcal{N}(0,\, 1) \text{ for all teams } k = 1,\dots,K 
\end{aligned}
\]

where $x^{\mathrm{vol}}_{j,t}$ is standardized $\log(1 + \text{FGA per 36\,min})$ for player $j$ in shot type $t$.

\vspace{1ex}

\textbf{Priors}
\[
\begin{aligned}
  \alpha_t &\sim \mathcal{N}(0,\; 0.25^2)
    & \beta_t  &\sim \mathcal{N}^{+}(1,\; 0.25^2) \\[4pt]
  \mu^{\mathrm{vol}}_t &\sim \mathcal{N}(0,\; 0.3^2) \\[4pt]
  \sigma^{\mathrm{p}}_t
    &\sim \mathcal{N}^{+}\!\bigl(0,\; s^{\mathrm{p}}_t\bigr),
    &s^{\mathrm{p}}
    &= (0.30,\; 0.30,\; 0.25) \\
  \sigma^{\mathrm{f}}_t
    % &\sim \mathcal{N}^{+}(0,\; 0.20) \;\; \forall\, t \\
    &\sim \mathcal{N}^{+}\!\bigl(0,\; s^{\mathrm{p}}_t\bigr),
    &s^{\mathrm{p}}
    &= (0.20,\; 0.20,\; 0.20) \\
  \sigma^{\mathrm{m}}_t
    &\sim \mathcal{N}^{+}\!\bigl(0,\; s^{\mathrm{m}}_t\bigr),
    &s^{\mathrm{m}}
    &= (0.12,\; 0.10,\; 0.06)
\end{aligned}
\]

% {\small Family-specific prior scales encode basketball domain knowledge: $s^{\mathrm{player}}$ and $s^{\mathrm{team}}$ vary by shot family because defensive scheme effects are largest at the rim and smallest on threes, while individual defender impact ($\sigma^{\mathrm{def}}_t$) uses a common scale across all families. The half-normal $\mathcal{N}^{+}$ constrains variance parameters to be positive. The slope prior $\beta_t \sim \mathcal{N}^{+}(1, 0.25^2)$ is centered at the identity transform, encoding the belief that a better Stage~1 prediction should correspond to a higher make probability.}

\end{block}

%----------------------------------------------------------------------------------------
% RESULTS
%----------------------------------------------------------------------------------------

\begin{block}{Results}

\textbf{Player rankings} with 90\% credible intervals (caterpillar plots):

\begin{figure}
\includegraphics[width=\linewidth]{player_caterpillar.png}
\caption{EPAA player effects by shot type with 90\% posterior intervals. Partial pooling shrinks low-volume players toward zero while preserving separation for high-volume stars.}
\end{figure}

\begin{figure}
\includegraphics[width=\linewidth]{defender_caterpillar.png}
\caption{ESAA defender effects by shot type with 90\% posterior intervals. Rim defense has greater variation than midrange and 3-pointers across players.}
\end{figure}

\textbf{Key findings:}
\begin{itemize}
\item EPAA re-ranks players substantially compared to raw eFG\%, rewarding efficient shooting on difficult looks and penalizing inflated stats from easy opportunities
\item Defender rankings reveal individual rim protection skill beyond team scheme effects
\item Volume--skill relationship is family-specific: jump shot volume correlates positively with skill (selection effect), while rim volume reflects role assignment rather than finishing talent
\end{itemize}

\end{block}

%----------------------------------------------------------------------------------------
% CONCLUSION / FUTURE WORK / CODE
%----------------------------------------------------------------------------------------

\begin{block}{Conclusion \& Future Work}

EPAA demonstrates that a two-stage framework---machine learning for context, Bayesian hierarchical modeling for talent---can \textbf{separate shot difficulty from shooting skill} with principled uncertainty quantification. The crossed random effects structure (shooter $\times$ defender $\times$ team) is a novel contribution to basketball analytics.

\vspace{0.5ex}

\textbf{Limitations \& future directions:}
\begin{itemize}
\item Posterior predictive intervals are slightly overconfident at high volumes; a game-level random intercept or beta-binomial likelihood is the principled fix
\item Split-half validation (first half $\to$ second half prediction) is underway as the headline stability metric
\item \textbf{True EPAA} extension: incorporating foul-drawing value via a two-tier architecture (foul probability + ``phantom shot'' modeling)
\end{itemize}

\end{block}

%----------------------------------------------------------------------------------------

\setbeamercolor{block alerted title}{fg=black,bg=norange}
\setbeamercolor{block alerted body}{fg=black,bg=white}

\begin{alertblock}{Code \& Reproducibility}

Full pipeline available on GitHub. Scan the QR code or visit:

\vspace{0.5ex}

\begin{center}
% Replace URL with your actual GitHub repo link
\qrcode[height=1.5in]{https://github.com/alex-susi/exp-fg-pct}

\vspace{0.3ex}
{\small \texttt{github.com/alex-susi/exp-fg-pct}}
\end{center}

\vspace{0.5ex}

{\small
\textbf{Pipeline:} R + Stan (\texttt{cmdstanr}) \\
\textbf{Scripts:} \texttt{01\_data\_pipeline.R} $\to$ \texttt{02\_xgb\_shot\_difficulty.R} $\to$ \texttt{03\_stan\_random\_effects.R} \\
\textbf{Data:} NBA Stats API via \texttt{hoopR}
}

\end{alertblock}

%----------------------------------------------------------------------------------------
% REFERENCES
%----------------------------------------------------------------------------------------

% \begin{block}{References}
% 
% {\small
% \begin{enumerate}
% \item Gelman, A., et al. (2013). \textit{Bayesian Data Analysis}, 3rd ed. CRC Press.
% \item Osborn, C. (2024). Probabilistic modeling of true hitting talent in MLB. \textit{CSAS}.
% \item Franks, A., et al. (2015). Characterizing the spatial structure of defensive skill in professional basketball. \textit{Annals of Applied Statistics}, 9(1), 94--121.
% \item Cervone, D., et al. (2016). A multiresolution stochastic process model for predicting basketball possession outcomes. \textit{JASA}, 111(514), 585--599.
% \end{enumerate}
% }
% 
% \end{block}

%----------------------------------------------------------------------------------------

\end{column}

\end{columns}

\end{frame}

\end{document}
